\section{Limitations}
\label{sec:limitations}

This is a model of a settlement's \emph{structure}, not a prediction of any
real one. The following limitations are material to how its results should be
read.

\paragraph{Parameters are chosen for legibility, not measured.}
Costs, capacities, growth rates and power draws were selected so that the
trade-offs between strategies are visible over a few hundred simulated days.
They are not derived from engineering estimates, and the absolute numbers in
Table~\ref{tab:econ} carry no physical meaning. What is meaningful is the
\emph{ordering} and the approximate ratios between strategies, since those
follow from the model's structure rather than from any single constant.

\paragraph{The physical grounding is qualitative.}
The model reproduces the qualitative structure of polar cold traps --- ice
where illumination never reaches, richer where it reaches least --- consistent
with observations of the lunar south polar region
\citep{paige2010diviner,colaprete2010water}. It does not attempt the actual
distribution, and its ice is far more abundant and more regular than any
measurement supports. Likewise, the case for subsurface habitation rests on
regolith shielding \citep{jia2010radiation,naito2020radiation} and on the
existence and stability of intact tubes
\citep{kaku2017intact,blair2017stability}, but the model represents neither
radiation dose nor structural stability at all. A tube
in this model cannot collapse, and its roof depth is recorded but unused. Any
statement here about the desirability of going underground is a statement about
\emph{land economics under the model's assumptions}, not about safety.

\paragraph{Illumination is static in time, but not in terrain.}
Solar exposure is an eight-direction average computed at generation and never
recomputed as the Sun moves; only the shading direction is animated. This is
defensible at polar latitudes, where the point is precisely that exposure
barely changes over a cycle, and progressively less defensible toward the
equator --- which is where the model's finding that mare sites cannot support
subsurface habitation is least trustworthy.

Exposure \emph{is} recomputed locally whenever the terrain changes, whether by
the player sculpting or by a meteor strike. Deposits, however, are seeded once
at generation and never re-seeded. The asymmetry is deliberate but has a
consequence worth stating: a player can manufacture shadow and cannot
manufacture ice, so excavating a crater produces ground that looks like an ice
site by every visible criterion and holds none. Whether that is the right
model of a cold trap is arguable --- ice accumulates in shadow over geological
time, not over a season --- but it is an assumption, not a derivation.

\paragraph{One tube per world.}
The model generates at most one lava tube, on one class of site, along the path
of a single rille. Real tube systems are networks. The strong result in
Section~\ref{sec:results} that tubes are polar-only is therefore a property of
this generator, not an inference about the Moon; what the model supports is the
narrower claim that a habitat strategy whose ceiling is set by geology behaves
differently from one whose ceiling is set by budget.

\paragraph{Provenance before instrumentation is lost.}
As set out in Section~\ref{sec:verification}, the number of simulation runs
performed during development was not recorded until late, and we decline to
reconstruct it.

\paragraph{The agent is one agent.}
The director results in Section~\ref{sec:failures} describe the behaviour of a
single heuristic policy. They are evidence about what that policy does, and
only suggestive about settlement dynamics in general.
